% Part: methods
% Chapter: induction
% Section: induction-on-N

\documentclass[../../../include/open-logic-section]{subfiles}

\begin{document}

\olfileid{mth}{ind}{inN}

\olsection{$\Nat$ 上的归纳法}

归纳法最简单的形式，是一种用来证明某个结果对所有自然数都成立的技巧。
它利用了这样一个事实：从 $0$ 开始反复加~$1$，我们最终能到达每一个自然数。
因此，要证明某个命题对每个数都成立，我们可以：(1)~证明它对 $0$ 成立；(2)~证明每当它对某个数~$n$ 成立时，它对下一个数~$n+1$ 也成立。
如果我们用 $P(n)$ 缩写“数~$n$ 具有性质 $P$”（并用 $P(k)$ 缩写“数~$k$ 具有性质 $P$”，等等），那么用归纳法证明“对所有 $n \in \Nat$，$P(n)$”就包括：
\begin{enumerate}
\item 证明 $P(0)$，以及
\item 证明，对任意~$k$，如果 $P(k)$ 则 $P(k+1)$。
\end{enumerate}

为了更清楚地说明这一点，假设我们同时有 (1) 和~(2)。
那么 (1) 告诉我们 $P(0)$ 为真。
如果我们还有~(2)，则特别地可知：若 $P(0)$ 则 $P(0+1)$，即 $P(1)$。
这是将一般性陈述“对任意~$k$，若 $P(k)$ 则 $P(k+1)$”中的~$k$ 取为~$0$ 而得到的。
因此，由假言推理，我们有~$P(1)$。
再次利用 (2)，这次将~$1$ 代入~$k$，我们得到：若 $P(1)$ 则~$P(2)$。
既然我们刚刚确立了~$P(1)$，由假言推理，我们就得到~$P(2)$。
依此类推。对于任意数~$n$，在这样操作 $n$~次之后，我们最终得到~$P(n)$。
所以，(1) 和~(2) 一起确立了 $P(n)$ 对任意 $n \in \Nat$ 成立。

让我们来看一个例子。假设我们想知道用 $n$ 个骰子可以掷出多少种不同的点数和。
尽管看起来有点傻，让我们从 $0$ 个骰子开始。
如果你没有骰子，那么你只能“掷”出一种可能的点数：完全没有点，总和为~$0$。
所以不同可能的掷出结果数是~$1$。
如果你只有一个骰子，即 $n=1$，有六种可能的值，从 $1$ 到~$6$。
用两个骰子，我们可以掷出 $2$ 到 $12$ 之间的任何总和，那是 $11$ 种可能。
用三个骰子，我们可以掷出 $3$ 到 $18$ 之间的任何数，即 $16$ 种不同的可能。
$1$，$6$，$11$，$16$：看起来像是有某种模式：也许答案是 $5n+1$？
当然，$5n+1$ 是最大的可能数，因为只有 $5n+1$ 个数介于 $n$（用 $n$ 个骰子能掷出的最小值，即全为 $1$）和 $6n$（能掷出的最大值，即全为 $6$）之间。

\begin{thm}
  用 $n$ 个骰子可以掷出 $n$ 与 $6n$ 之间的所有 $5n+1$ 个可能的值。
\end{thm}

\begin{proof}
令 $P(n)$ 为断言：“使用 $n$~个骰子可以掷出 $n$ 与~$6n$ 之间的任何数。”
为了使用归纳法，我们证明：
\begin{enumerate}
\item \emph{归纳基始} $P(1)$，即仅用一个骰子，你可以掷出 $1$ 到 $6$ 之间的任何数。
\item \emph{归纳步骤}，对于所有 $k$，若 $P(k)$ 则~$P(k+1)$。
\end{enumerate}

(1) 通过观察一个六面骰子即可证明。它有六个面，并且 $1$ 到~$6$ 之间的每个数都出现在某个面上。因此，使用单个骰子可以掷出 $1$ 到~$6$ 之间的任何数。

为了证明 (2)，我们假设条件句的前件，即 $P(k)$。这个假设称为\emph{归纳假设}。我们用它来证明~$P(k+1)$。困难的部分在于，找到一种方法，用掷 $k$~个骰子的可能值加上额外的第 $k+1$ 个骰子掷出的结果，来思考掷 $k+1$ 个骰子的可能值——不过，如果想使用归纳假设，我们就必须这么做。

归纳假设表明，我们可以用 $k$~个骰子得到 $k$ 到~$6k$ 之间的任何数。如果我们的第 $(k+1)$ 个骰子掷出了~$1$，这会使总和增加~$1$。因此，我们可以先掷 $k$~个骰子，然后在第 $(k+1)$ 个骰子上掷出~$1$，从而掷出 $k+1$ 到 $6k+1$ 之间的任何值。
剩下的是什么值呢？是 $6k+2$ 到 $6k+6$ 这些值。我们可以通过在 $k$ 个骰子上都掷出 $6$，然后在第 $(k+1)$ 个骰子上掷出一个 $2$ 到 $6$ 之间的数来得到这些值。
综合起来，这意味着用 $k+1$ 个骰子，我们可以掷出 $k+1$ 到 $6(k+1)$ 之间的任何数，也就是说，我们在归纳假设~$P(k)$ 的前提下，证明了~$P(k+1)$。
\end{proof}

我们经常在想要证明关于一系列对象（数、集合等）的某个性质时使用归纳法，而这些对象本身是“归纳地”定义的，即用第 $n$ 个对象来定义第 $(n+1)$ 个对象。
例如，我们可以用以下方式定义到~$n$ 为止的自然数之和~$s_n$：
\begin{align*}
  s_0 & = 0\\
  s_{n+1} & = s_n + (n+1)
\end{align*}
这个定义给出：
\begin{align*}
  s_0 & = 0,\\
  s_1 & = s_0 + 1 && = 1,\\
  s_2 & = s_1 + 2 && = 1 + 2 = 3\\
  s_3 & = s_2 + 3 && = 1 + 2 + 3 = 6, \text{ 等等。}
\end{align*}
现在我们可以用归纳法证明，$s_n = n(n+1)/2$。

\begin{prop}
  $s_n = n(n+1)/2$。
\end{prop}

\begin{proof}
  我们需要证明（1）$s_0 = 0\cdot(0 + 1)/2$，以及（2）如果 $s_k =
  k(k+1)/2$，那么 $s_{k+1} = (k+1)(k+2)/2$。（1）是显然的。要证明（2），我们假定归纳假设成立：$s_k = k(k+1)/2$。利用它，我们需要证明 $s_{k+1} = (k+1)(k+2)/2$。

  $s_{k+1}$ 是什么呢？根据定义，$s_{k+1} = s_k + (k+1)$。由归纳假设，$s_k = k(k+1)/2$。我们可以将其代入前一个等式，然后只需稍作分数的算术运算：
  \begin{align*}
    s_{k+1} & = \frac{k(k+1)}{2} + (k+1) = {}\\
    & = \frac{k(k+1)}{2} + \frac{2(k+1)}{2} = {}\\
    & = \frac{k(k+1) + 2(k+1)}{2} = {}\\
    & = \frac{(k+2)(k+1)}{2}.
  \end{align*}
\end{proof}

这里的重要经验是：如果你要证明关于某个归纳定义的序列 $a_n$ 的性质，归纳法显然是首选。即使序列不是归纳定义的（例如骰子点数的可能性问题），只要你能以某种方式将第~$k+1$ 种情况与第~$k$ 种情况联系起来，你仍然可以使用归纳法。

\end{document}